\chapter{Assessing Risks and Limitations of Kernel-Level Anti-Cheats
using Windows-Kernel-Explorer}
\label{chap:two}

\epigraph{\itshape The programmers of tomorrow are the wizards of the future. You're going to look like you have magic powers compared to everybody else.}{-- Gabe NEWELL}

\stopcontents[chapters]

\startcontents[chapters]

\printMyMinitoc

\fancyhead[LE]{\textsc{\chaptername~\thechapter}}

\section{Structural Limitations of Kernel-Level Anti-Cheats}
\label{sec:two:one}

\subsection{Technological Arms Race}
\label{sec:two:one:one}

The deployment of kernel-level anti-cheat systems did not resolve the cheating problem: it escalated it. By elevating the detection layer to Ring 0, publishers effectively forced cheat developers to match that privilege level, triggering a continuous cycle of technical escalation in which each defensive advancement is systematically studied, reverse-engineered, and circumvented. This structural escalation is not merely theoretical: systematic grey-box testing of leading anti-cheat systems shows that even kernel-level solutions remain bypassed by a range of techniques, and that VAC (the system protecting CS2) ranks among the weakest evaluated, detecting only a narrow subset of the tested attack vectors \cite{Collins2024}.

This dynamic is structurally inherent to the paradigm. Once deployed, an anti-cheat driver enforces a static set of invariants (known cheat signatures, protected memory regions, hooked system calls, integrity checks) fixed at the moment of release. Cheat developers, operating in a well-funded and motivated ecosystem, treat each new version as a reverse-engineering target, often using the same kernel-level tooling described in Section 1.1.3 to map the detection surface before it can be updated.

The BYOVD technique (Section 1.1.3) is a direct product of this arms race: it emerged once anti-cheats began enforcing driver signature validation, prompting cheat developers to exploit the pre-existing trust chain of legitimately signed third-party drivers. Vendors responded with blocklists of known-vulnerable drivers — a measure Microsoft has since integrated at the \gls{OS} level via its Vulnerable Driver Blocklist — while cheat developers shifted to newer, uncatalogued vulnerable drivers. Each iteration demands new engineering investment from the anti-cheat side, while the cheat ecosystem distributes evasion research across a community sharing tools openly. This asymmetry is documented in practice: cheat sellers maintain live detection-rate status pages, and a 37-day monitoring period shows at least one cheat remaining operational 100\% of the time against kernel-protected titles, with median cheat uptime across all studied games above 85\% \cite{Collins2024}.

This asymmetry reflects a deeper structural weakness: an anti-cheat update requires rigorous stability validation, since a defective kernel component can trigger system-wide crashes, corrupt system state, or introduce new vulnerabilities. The cheat ecosystem faces no such constraint: a new evasion technique can ship within days, while the corresponding patch may take weeks or months to safely deploy. The defender is therefore structurally slower than the attacker, regardless of resources invested, and each new game release, OS update, or hardware generation compounds this dependency: an unmaintained kernel component quickly becomes a liability, its Ring 0 presence introducing risk without corresponding detection benefit.

\subsection{Limits Against New Forms of Cheating}
\label{sec:two:one:two}

Beyond the arms race, kernel-level anti-cheats are architecturally incapable of addressing cheat vectors that externalize execution outside the host OS entirely: a limitation rooted not in implementation quality but in the threat model itself.

Kernel-level anti-cheats assume the cheat leaves a detectable footprint within the host OS: unauthorized memory modification, suspicious process activity, unsigned driver loading, or anormal system function calls. This assumption breaks down entirely against \gls{DMA} cheats (Section 1.1.3, Figure 5). A DMA cheat connects a dedicated PCIe card directly to the host's memory bus, reading process memory (including the decrypted Entity List) at the hardware level, bypassing the CPU and OS process model. Game state is transmitted to a physically separate machine for overlay rendering and aim computation, with synthetic inputs fed back via a hardware spoofer. Since no software executes on the host, no driver loads, no process is injected, and no local memory is modified, the kernel-level anti-cheat has no artifact to detect, since its Ring 0 monitoring simply never reaches the hardware interface involved.

This is a categorical defeat against a commercially available and increasingly prevalent class of cheats: a functional DMA setup (PCIe card, secondary machine, input spoofer) costs a few hundred dollars, well within reach of motivated cheaters in games where ranked progression carries economic or reputational stakes. Software-side countermeasures (monitoring PCIe enumeration, measuring access-latency anomalies, detecting known DMA firmware signatures) share the same fundamental weakness as the broader arms race: they are signature-based, enumerable, and bypassable, since the attacker controls the hardware itself.

The architectural conclusion follows directly: kernel-level anti-cheat can only monitor the host OS processes. Any vector operating below that boundary (hardware memory access, external processing) falls entirely outside even the most privileged driver's detection perimeter.

\subsection{Economic and Operational Costs}
\label{sec:two:one:three}

The structural limitations described above do not manifest in a vacuum: they translate into concrete and ongoing economic costs for publishers who choose to deploy kernel-level anti-cheat systems. These costs operate across three distinct dimensions that collectively challenge the long-term viability of the approach.

The first is driver maintenance. A kernel driver is not a product that can be shipped and forgotten: it must be continuously adapted to new OS versions, new hardware generations, and new security policies enforced by Microsoft's evolving Secure Boot and driver signing infrastructure. Each Windows update cycle introduces potential compatibility breaks that must be diagnosed and patched before they manifest as crashes or silent detection failures in the player base. This maintenance burden requires a dedicated engineering team with specialized kernel-level expertise — a skillset that is both scarce and expensive in the current software labor market.

The second is security incident management. As discussed in Appendix D.3, the mandatory installation of a Ring 0 component exposes the publisher to significant liability in the event of a security flaw. The historical record is unambiguous: the Capcom anti-cheat driver exploit allowed arbitrary unsigned code execution at kernel level on every machine where it was installed; the Genshin Impact anti-cheat driver was weaponized by a third party to deploy ransomware by exploiting its legitimate Ring 0 access. Each such incident triggers a crisis response cycle — emergency patching, public disclosure management, potential regulatory scrutiny, and, in severe cases, class-action litigation. The cost of a single major incident can exceed the cumulative engineering investment in the anti-cheat system itself.

The third is reputational and commercial impact. As detailed in Appendix D.2, a growing segment of the player base actively rejects kernel-level solutions, translating publisher anti-cheat decisions into measurable adoption losses. The exclusion of Linux and Steam Deck users (a platform whose market share has grown substantially with the mainstream success of Valve's handheld) from titles enforcing kernel drivers represents a quantifiable reduction in the addressable player pool. This exclusion is compounded by the architectural characteristics of the most aggressive solutions: always-on kernel drivers that boot independently of the game process and that, in the case of FACEIT Anti-Cheat, require users to disable Windows memory integrity protections and Hyper-V — a condition that exposes the system to a broader class of threats than those it is intended to prevent \cite{Dorner2024}. For free-to-play titles where player count directly drives monetization through network effects and cosmetic economies, this trade-off carries direct financial consequences that compound over time.

Taken together, these three cost vectors (maintenance, incident response, and adoption friction) establish that kernel-level anti-cheat is not simply a technical solution with technical limitations, but a strategic commitment with sustained organizational costs. The central question for publishers is therefore whether the detection effectiveness gained justifies this investment, particularly given the categorical limitations against hardware-assisted cheating vectors identified in Section 2.1.2. It is precisely this question that motivates the server-side behavioral approach explored in Chapter 3.

\section{Experimental Section: Illustrating Potential Intrusiveness}
\label{sec:two:two}

\subsection{Experimental Objective}
\label{sec:two:two:one}

\begin{itemize}
\item Illustrate the level of access of a privileged component
\item Demonstrate security and privacy implications
\item Headstart thanks to https://github.com/AxtMueller/Windows-Kernel-Explorer --> PROBLEM: ONLY SUPPORTED UNTIL WINDOWS 10. WHAT DO DO INSTEAD? I SEARCHED FOR ALTERNATIVES (READ RAM CONTENT ON LINUX, OTHER TOOLS TO READ RAM CONTENT ON WINDOWS) BUT COULDN'T GET ANYTHING TO WORK ON LINUX NOR WINDOWS. I COULD SHOW HOW CHEAT ENGINE CAN FIND A VARIABLE LOCATION IN RAM FROM USING SUCCESIVE VALUES TO PINPOINT THE EXACT LOCATION OF THE VARIABLE IN RAM (MEANING YOU CAN THEN READ IT OR MODIFY IT) BUT WOULD THAT BE REALLY USEFULL IN THE CONTENT OF THIS THESIS?
\end{itemize}

\subsection{Experimental Observations}
\label{sec:two:two:two}

\begin{itemize}
\item Access to sensitive system resources
\item Privilege differences compared to regular software
\item Observation of non-game-related data
\end{itemize}

\subsection{Critical Analysis}
\label{sec:two:two:three}

\begin{itemize}
\item Real vs theoretical implications
\item Link with user perception
\item Discussion of effectiveness vs intrusiveness proportionality
\end{itemize}
